\chapter{Harmonic Forms and the Hodge Theorem}\label{hodge-theorem}
\chapterstatus{complete}

\section{Introduction}\label{hodge-theorem:section-introduction}

On a compact oriented Riemannian manifold every de Rham cohomology class contains exactly one harmonic form
(Theorem~\ref{hodge-theorem:theorem-hodge}). This is the theorem that makes Hodge theory possible: it turns a topological object,
the cohomology of a manifold, into the solution space of an elliptic equation, on which the metric and, in the complex case, the
complex structure act. The consequences drawn here are the Hodge decomposition of the space of forms, Poincar\'e duality via the
Hodge star, and Bochner's vanishing theorem; the complex analogue is in Chapter~\ref{hodge-decomposition}. References:
\cite[Chapter 6]{warner-foundations}, \cite[Chapter 4]{wells-differential-analysis}, \cite[Chapter 0]{griffiths-harris},
\cite[Chapter 5]{voisin-hodge-1}.

\noindent\textbf{Prerequisites.} Chapter~\ref{elliptic-operators} (Elliptic Operators on Compact Manifolds) and
Chapter~\ref{de-rham} (de Rham Cohomology).

Throughout, $(M, g)$ is a compact oriented Riemannian $n$-manifold, with the $L^2$ inner product, the codifferential $\delta$ and the
Laplacian $\Delta = d\delta + \delta d$ of Definition~\ref{riemannian-geometry:definition-laplacian}, and
$\mathcal{H}^k = \mathcal{H}^k(M, g) = \{\alpha \in \Omega^k(M) : \Delta\alpha = 0\}$.

\section{The Hodge decomposition}\label{hodge-theorem:section-decomposition}

\begin{theorem}\label{hodge-theorem:theorem-decomposition}
For each $k$ the space $\mathcal{H}^k$ is finite-dimensional and consists of the forms with $d\alpha = 0$ and $\delta\alpha = 0$, and there is an orthogonal direct sum decomposition
\[
\Omega^k(M) = \mathcal{H}^k \oplus d\Omega^{k-1}(M) \oplus \delta\Omega^{k+1}(M),
\]
with $\ker d = \mathcal{H}^k \oplus d\Omega^{k-1}$ and $\ker\delta = \mathcal{H}^k \oplus \delta\Omega^{k+1}$. Moreover there is a unique operator $G \colon \Omega^k(M) \to \Omega^k(M)$, the \emph{Green operator},
that vanishes on $\mathcal{H}^k$, maps $(\mathcal{H}^k)^\perp$ into itself, and satisfies $\alpha = H\alpha + \Delta G\alpha$, where $H$ is the orthogonal projection onto $\mathcal{H}^k$; it commutes with $d$,
$\delta$ and $\Delta$.
\end{theorem}

\begin{proof}
$\Delta$ is elliptic, formally self-adjoint and non-negative (Example~\ref{elliptic-operators:example-elliptic} and
Proposition~\ref{riemannian-geometry:proposition-adjoint}), so by Theorem~\ref{elliptic-operators:theorem-fredholm} its kernel $\mathcal{H}^k$ is finite-dimensional and consists of smooth forms, and
\[
\Omega^k(M) = \mathcal{H}^k \oplus \Delta\Omega^k(M)
\]
orthogonally: indeed $L^2 = \operatorname{im}\Delta \oplus \ker\Delta$ by Theorem~\ref{elliptic-operators:theorem-fredholm}(3), and a smooth form orthogonal to $\mathcal{H}^k$ is $\Delta\beta$ for a smooth $\beta$
by elliptic regularity. The description of $\mathcal{H}^k$ as $\ker d \cap \ker\delta$ is
Proposition~\ref{riemannian-geometry:proposition-adjoint}(3).

Since $\Delta\beta = d(\delta\beta) + \delta(d\beta)$, the second summand is contained in $d\Omega^{k-1} + \delta\Omega^{k+1}$, so the three spaces span $\Omega^k$. They are pairwise orthogonal: harmonic forms are
orthogonal to $d\Omega^{k-1}$ because $(\alpha, d\beta) = (\delta\alpha, \beta) = 0$, and similarly to $\delta\Omega^{k+1}$; and $(d\beta, \delta\gamma) = (dd\beta, \gamma) = 0$. Hence the sum is direct and
orthogonal. If $d\alpha = 0$ and $\alpha = \eta + d\beta + \delta\gamma$, then $0 = (d\alpha, \gamma) = (d\delta\gamma, \gamma) = \|\delta\gamma\|^2$, so $\alpha \in \mathcal{H}^k \oplus d\Omega^{k-1}$; the statement for
$\ker\delta$ is symmetric.

For the Green operator, define $G\alpha = 0$ on $\mathcal{H}^k$ and, on the orthogonal complement, $G\alpha = \beta$ where $\beta$ is the unique solution of $\Delta\beta = \alpha$ orthogonal to $\mathcal{H}^k$
(Theorem~\ref{elliptic-operators:theorem-fredholm}(3) gives existence and smoothness, and uniqueness holds because two solutions differ by a harmonic form). Then $\alpha = H\alpha + \Delta G\alpha$ by
construction. Since $\Delta$ commutes with $d$, $\delta$ and $\Delta$, and these operators preserve $\mathcal{H}^\bullet$ and its orthogonal complement (for instance $d\mathcal{H}^k = 0$ and
$(d\alpha, \eta) = (\alpha, \delta\eta) = 0$ for $\eta$ harmonic), the uniqueness of $G$ forces $Gd = dG$, $G\delta = \delta G$ and $G\Delta = \Delta G$.
\end{proof}

\begin{theorem}[Hodge]\label{hodge-theorem:theorem-hodge}
The map $\mathcal{H}^k \to H^k_{\mathrm{dR}}(M)$, $\alpha \mapsto [\alpha]$, is an isomorphism: every de Rham class contains a unique harmonic representative, namely the one of smallest $L^2$ norm. In
particular $H^k_{\mathrm{dR}}(M)$ is finite-dimensional, and $\dim\mathcal{H}^k = b_k(M)$.
\end{theorem}

\begin{proof}
Harmonic forms are closed, so the map is defined. It is injective: if $\alpha \in \mathcal{H}^k$ is exact, $\alpha = d\beta$, then $\|\alpha\|^2 = (\alpha, d\beta) = (\delta\alpha, \beta) = 0$. It is surjective: a closed
$\alpha$ decomposes as $\alpha = \eta + d\beta + \delta\gamma$ with $\eta$ harmonic, and $\delta\gamma = 0$ as in the proof of Theorem~\ref{hodge-theorem:theorem-decomposition}, so $\alpha = \eta + d\beta$ is
cohomologous to $\eta$. If $\alpha = \eta + d\beta$ is any representative, then $\|\alpha\|^2 = \|\eta\|^2 + \|d\beta\|^2 \geq \|\eta\|^2$ with equality only for $\alpha = \eta$, which is the statement about the
norm. Finite-dimensionality and the identification with the Betti numbers follow from Theorem~\ref{de-rham:theorem-de-rham}.
\end{proof}

\begin{corollary}\label{hodge-theorem:corollary-poincare-duality}
The Hodge star induces an isomorphism $* \colon \mathcal{H}^k \to \mathcal{H}^{n-k}$, hence $b_k(M) = b_{n-k}(M)$, and the pairing
\[
H^k_{\mathrm{dR}}(M) \times H^{n-k}_{\mathrm{dR}}(M) \to \RR, \qquad ([\alpha], [\beta]) \mapsto \int_M\alpha \wedge \beta,
\]
is nondegenerate; for $[\beta] = [*\alpha]$ with $\alpha$ harmonic and nonzero it takes the value $\|\alpha\|^2 > 0$.
\end{corollary}

\begin{proof}
$*$ commutes with $\Delta$ (Proposition~\ref{riemannian-geometry:proposition-adjoint}(4)) and is an isomorphism of $\Omega^k$ onto $\Omega^{n-k}$, hence restricts to an isomorphism on harmonic forms; combine
with Theorem~\ref{hodge-theorem:theorem-hodge}. For a nonzero harmonic $\alpha$, $\int_M\alpha \wedge *\alpha = \|\alpha\|^2 > 0$, so the pairing is nondegenerate on the harmonic representatives, which represent
all classes.
\end{proof}

\begin{remark}\label{hodge-theorem:remark-dependence-on-metric}
The space $\mathcal{H}^k$ depends on the metric, but its dimension does not, and the isomorphism of Theorem~\ref{hodge-theorem:theorem-hodge} identifies all of them with $H^k_{\mathrm{dR}}(M)$. This is the
mechanism that produces the Hodge structures of Chapter~\ref{hodge-decomposition}: on a compact K\"ahler manifold the harmonic forms decompose by type, and the resulting decomposition of cohomology is
independent of the metric, although the harmonic representatives are not.
\end{remark}

\begin{proposition}\label{hodge-theorem:proposition-lie-group-harmonic}
Let $G$ be a compact connected Lie group with a bi-invariant Riemannian metric (Example~\ref{riemannian-geometry:example-metrics}) and an orientation. Then the harmonic forms on $G$ are exactly the bi-invariant forms.
\end{proposition}

\begin{proof}
Left and right translations are isometries, and they preserve the orientation, being connected to the identity through diffeomorphisms; so they commute with the Hodge star, and $*$ maps bi-invariant forms to bi-invariant forms. A bi-invariant
form $\omega$ is closed (Theorem~\ref{de-rham:theorem-lie-group-cohomology}), and so is $*\omega$; hence $\delta\omega = \pm{*}d{*}\omega = 0$ and $\omega$ is harmonic (Proposition~\ref{riemannian-geometry:proposition-adjoint}(3)). The
bi-invariant $k$-forms and the harmonic $k$-forms both map isomorphically onto $H^k_{\mathrm{dR}}(G)$ (Theorems~\ref{de-rham:theorem-lie-group-cohomology} and~\ref{hodge-theorem:theorem-hodge}), so the inclusion of the first
space into the second is an equality.
\end{proof}

\begin{counterexample}\label{hodge-theorem:counterexample-compactness}
Compactness is essential in the Hodge theorem. On $\RR$ with the Euclidean metric, the harmonic functions are the affine functions $ax + b$, a two-dimensional space, while $H^0_{\mathrm{dR}}(\RR) = \RR$; and the harmonic $1$-forms $(ax + b)\,dx$
form a two-dimensional space although $H^1_{\mathrm{dR}}(\RR) = 0$. Without compactness harmonic forms need not be closed or coclosed: the function $x$ is harmonic with $dx \neq 0$, and $x\,dx$ is harmonic with $\delta(x\,dx) = -1 \neq 0$, since $\delta = -{*}d{*}$ on $1$-forms on $\RR$. The
integration by parts behind Proposition~\ref{riemannian-geometry:proposition-adjoint}(3) fails for forms without compact support.
\end{counterexample}

\section{Applications}\label{hodge-theorem:section-applications}

\begin{corollary}\label{hodge-theorem:corollary-euler-characteristic}
For a compact oriented manifold, $\sum_k(-1)^k\dim\mathcal{H}^k = \chi(M)$, which is the index of the elliptic operator
$d + \delta \colon \Omega^{\mathrm{even}}(M) \to \Omega^{\mathrm{odd}}(M)$.
\end{corollary}

\begin{proof}
The first equality is Theorem~\ref{hodge-theorem:theorem-hodge} together with the definition of $\chi$ as the alternating sum of Betti numbers
(Corollary~\ref{poincare-duality:corollary-euler-odd}). For the second, $(d + \delta)^* = d + \delta$ exchanges even and odd forms, its kernel on even forms is $\bigoplus_{k \text{ even}}\mathcal{H}^k$ and on odd
forms $\bigoplus_{k \text{ odd}}\mathcal{H}^k$, because $(d + \delta)\alpha = 0$ implies $\Delta\alpha = (d+\delta)^2\alpha = 0$; so the index is $\sum_k(-1)^k\dim\mathcal{H}^k$.
\end{proof}

\begin{theorem}[Bochner]\label{hodge-theorem:theorem-bochner}
Let $(M, g)$ be a compact oriented Riemannian manifold with $\mathrm{Ric} \geq 0$. Then every harmonic $1$-form is parallel, and $b_1(M) \leq n$. If moreover $\mathrm{Ric} > 0$ at some point, then
$b_1(M) = 0$.
\end{theorem}

\begin{proof}
By the Weitzenb\"ock formula (Remark~\ref{riemannian-geometry:remark-weitzenbock}), on $1$-forms $\Delta = \nabla^*\nabla + \mathrm{Ric}$. For $\alpha$ harmonic,
\[
0 = (\Delta\alpha, \alpha) = \|\nabla\alpha\|^2 + \int_M\mathrm{Ric}(\alpha^\sharp, \alpha^\sharp)\,\mathrm{vol},
\]
where $\alpha^\sharp$ is the vector field dual to $\alpha$. Both terms are non-negative, hence both vanish: $\nabla\alpha = 0$, so $\alpha$ is parallel and determined by its value at one point, giving
$\dim\mathcal{H}^1 \leq n$; and $\mathrm{Ric}(\alpha^\sharp, \alpha^\sharp) = 0$ everywhere. If $\mathrm{Ric} > 0$ at $p$ and $\alpha \neq 0$, then $\alpha$ is nowhere zero (being parallel and nonzero somewhere), so
$\mathrm{Ric}(\alpha^\sharp, \alpha^\sharp)(p) > 0$, a contradiction; hence $\mathcal{H}^1 = 0$ and $b_1(M) = 0$ by Theorem~\ref{hodge-theorem:theorem-hodge}.
\end{proof}

\begin{example}\label{hodge-theorem:example-applications}
\begin{enumerate}
\item On the round sphere $S^n$ with $n \geq 2$, $\mathrm{Ric} > 0$, so $b_1 = 0$, which agrees with Example~\ref{de-rham:example-spheres}.
\item On a flat torus the harmonic forms are those with constant coefficients (Example~\ref{riemannian-geometry:example-harmonic-forms}), so Theorem~\ref{hodge-theorem:theorem-hodge} recovers
$b_k(T^n) = \binom{n}{k}$.
\item A compact oriented manifold with a metric of positive Ricci curvature has $b_1 = 0$; by Bochner's argument applied to the other Laplacians, positivity assumptions on curvature produce vanishing
theorems in all degrees. In the complex setting the same mechanism gives the Kodaira vanishing theorem (Chapter~\ref{kodaira}).
\end{enumerate}
\end{example}

\begin{remark}\label{hodge-theorem:remark-elliptic-complexes}
The proof of Theorem~\ref{hodge-theorem:theorem-decomposition} used only that $(\Omega^\bullet, d)$ is a complex of sections of vector bundles whose Laplacian is elliptic: an \emph{elliptic complex}. The same
argument applies verbatim to any elliptic complex, in particular to the Dolbeault complex $(\Omega^{p,\bullet}, \bar\partial)$ on a compact complex manifold, giving the isomorphism between Dolbeault
cohomology and $\bar\partial$-harmonic forms (Chapter~\ref{dolbeault}), which is how the Hodge decomposition of Chapter~\ref{hodge-decomposition} is proved. There is also a proof of the Hodge theorem by the
heat equation, which constructs the harmonic projection as $\lim_{t \to \infty}e^{-t\Delta}$ and yields the local index formula \cite{lawson-michelsohn}.
\end{remark}

\section{Exercises}\label{hodge-theorem:section-exercises}

\begin{exercise}\label{hodge-theorem:exercise-harmonic-top}
Show that on a compact connected oriented Riemannian $n$-manifold, $\mathcal{H}^0 = \RR$ and $\mathcal{H}^n = \RR\,\mathrm{vol}_g$, and deduce $b_0 = b_n = 1$.
\end{exercise}

\begin{solution}
A harmonic function has $df = 0$, hence is constant; apply $*$ for the top degree (Corollary~\ref{hodge-theorem:corollary-poincare-duality}).
\end{solution}

\begin{exercise}\label{hodge-theorem:exercise-green-properties}
Show that the Green operator is self-adjoint, that it is compact as an operator on $L^2$, and that $\alpha \mapsto H\alpha$ is the orthogonal projection onto $\mathcal{H}^k$ with $H = I - \Delta G$.
\end{exercise}

\begin{exercise}\label{hodge-theorem:exercise-harmonic-not-metric-independent}
Give an example of a compact manifold and two metrics for which the harmonic representatives of a fixed cohomology class differ.
\end{exercise}

\begin{solution}
On a torus $\RR^2/\ZZ^2$ with the flat metric, the harmonic representative of the class of $dx$ is $dx$; after a conformal change $g' = e^{2f}g$ with non-constant $f$, the codifferential changes and the
harmonic representative of the same class is a different $1$-form, since $\delta'$ differs from $\delta$.
\end{solution}

\begin{exercise}\label{hodge-theorem:exercise-hodge-riemann-surface}
Let $M$ be a compact orientable surface of genus $g$. Using Theorem~\ref{hodge-theorem:theorem-hodge} and
Corollary~\ref{hodge-theorem:corollary-poincare-duality}, show that $\dim\mathcal{H}^1 = 2g$ and that the intersection pairing on $\mathcal{H}^1$ is a nondegenerate alternating form.
\end{exercise}
